\documentclass[11pt,a4paper]{article}
\usepackage{fontspec,amsmath,amssymb,amsthm,microtype}
\setmainfont{lmroman10-regular.otf}[BoldFont=lmroman10-bold.otf,ItalicFont=lmroman10-italic.otf,BoldItalicFont=lmroman10-bolditalic.otf]
\usepackage[margin=25mm]{geometry}
\usepackage[unicode,hidelinks]{hyperref}
\hypersetup{pdfauthor={LangevinLean Project}}
\setlength{\emergencystretch}{3em}
\setlength{\parskip}{0.4em}
\setlength{\parindent}{0pt}
\allowdisplaybreaks[2]
\raggedbottom
\title{\texorpdfstring{Unadjusted Langevin under Strong Convexity}{Unadjusted Langevin under Strong Convexity}}
\author{LangevinLean Project}
\date{}
\begin{document}\maketitle
Draft mathematical exposition; these notes may contain mistakes.
The strongly convex convergence theorems construct their Brownian reference and prove Gibbs stationarity. The LSI extension derives its entropy bound from smoothness, LSI, and finite initial KL.
\section{\texorpdfstring{Theorem}{Theorem}}


Let \(d\ge1\), and let \(U:\mathbb R^d\to\mathbb R\) be differentiable and \(m\)-strongly convex, with \(m>0\) and \(L\)-Lipschitz gradient. These assumptions imply \(m\le L\). Define


\[\pi(dx)=Z^{-1}e^{-U(x)}dx,\qquad Z=\int e^{-U(x)}dx\in(0,\infty).\]


Let \(\mu_0\) be a probability law with \(W_2(\mu_0,\pi)<\infty\). A finite second moment suffices. For \(h>0\), take independent standard Gaussian vectors \(\xi_k\), independent of initialization, and set


\[X_{k+1}=X_k-h\nabla U(X_k)+\sqrt{2h}\,\xi_k,\qquad\mu_k=\operatorname{Law}(X_k).\]


The transition kernel is \(P_h(x,\cdot)=\mathcal N(x-h\nabla U(x),2hI_d)\), so \(\mu_{k+1}=\mu_kP_h\). The independence assumptions ensure that conditioning on \(X_k\) gives exactly this kernel. Thus a bound for the iterated probability laws is also a bound for the displayed stochastic recursion.




If \(0<h\le2/(m+L)\), then for every \(k\in\mathbb N\),

\[W_2(\mu_k,\pi)\le(1-mh)^kW_2(\mu_0,\pi)+\frac{33}{20}\frac Lm\sqrt{hd}.\]

The rational coefficient \(33/20\) is exactly \(1.65\). The initial term uses the true Wasserstein infimum over couplings.


\section{\texorpdfstring{From convexity and smoothness to sharp contraction}{From convexity and smoothness to sharp contraction}}



Write \(g=\nabla U\). Strong convexity and gradient Lipschitzness give the complementary quadratic models


\[U(x)+\langle g(x),y-x\rangle+\frac m2\|y-x\|^2
\le U(y)\le U(x)+\langle g(x),y-x\rangle+\frac L2\|y-x\|^2.\]


The lower inequality follows by subtracting \(m\|x\|^2/2\) and using the supporting-plane inequality for a differentiable convex function. For the upper inequality, integrate the gradient along the segment from \(x\) to \(y\); the remainder is bounded by \(L\int_0^1t\,dt\,\|y-x\|^2\). Thus no Hessian assumption is required.




Apply cocoercivity to the shifted function \(U-m\|\cdot\|^2/2\), whose quadratic upper-model constant is \(L-m\). With \(d=x-y\) and \(v=g(x)-g(y)\), expanding its conclusion yields


\[mL\|d\|^2+\|v\|^2\le(m+L)\langle v,d\rangle.\]


Completing the square makes the geometry explicit:


\[\left\|v-\frac{m+L}{2}d\right\|\le\frac{L-m}{2}\|d\|.\]


Consequently the triangle inequality bounds \(\|d-hv\|\) by


\[\left(\left|1-\frac{m+L}{2}h\right|+\frac{L-m}{2}h\right)\|d\|.\]


For \(0\le h\le2/(m+L)\), the factor is \(1-mh\). For \(2/(m+L)\le h\le2/L\), it is \(Lh-1\). The boundary \(m=L\) is included because the centered ball then has radius zero; there is no division by \(L-m\).



\section{\texorpdfstring{Normalize the target and control its score}{Normalize the target and control its score}}



Strong convexity gives, at any fixed \(x_0\), a quadratic lower bound for \(U(x)\). Completing its square shows that \(e^{-U(x)}\), even when multiplied by any fixed quadratic polynomial in \(x\), is dominated by an integrable translated Gaussian. Therefore the partition function is finite, strictly positive, and the normalized Gibbs law has a finite second moment (see \cite[Section~3.26.13]{LangevinLean2026ULASC} and \cite[Section~3.26.14]{LangevinLean2026ULASC}).




A second moment needed below concerns the gradient. The smooth upper model implies, for any \(v\),


\[U(x-v)\le U(x)-\langle g(x),v\rangle+\frac L2\|v\|^2.\]


Exponentiate the rearranged inequality and integrate. Translation of Lebesgue measure preserves the partition function, giving


\[\int e^{\langle g(x),v\rangle}\,\pi(dx)\le e^{L\|v\|^2/2}.\]


For each unit coordinate vector, use both signs of a scalar multiple of that vector. The resulting sub-Gaussian bound yields a centered directional score with second moment at most \(L\). Summing over an orthonormal basis gives


\[\int\|g(x)\|^2\,\pi(dx)\le Ld.\]


This derivation uses the original first-derivative assumptions. It also proves the gradient integrability that will justify drift integrals, rather than treating that integrability as an additional assumption.



\section{\texorpdfstring{Construct a stationary comparison diffusion}{Construct a stationary comparison diffusion}}



Take the constructed centered Gaussian process \(B_t\), with independent increments, covariance \(\mathbb E[B_sB_t^{\mathsf T}]=\min(s,t)I_d\), and continuous dependence on time in mean square. Choose it independently of \(Y_0\sim\pi\). In the Hilbert space of square-integrable random vectors, solve


\[Y_t=Y_0-\int_0^t g(Y_s)\,ds+\sqrt2 B_t.\]


The drift induces an \(L\)-Lipschitz map on that Hilbert space. Picard iteration converges on finite time intervals, and the integral inequality for the difference of two candidate solutions gives uniqueness by Gronwall. The resulting process is continuous in mean square. Brownian increments on disjoint intervals are independent and have covariance equal to the interval length.




Stationarity needs a separate argument. Start a Metropolis-adjusted Euler chain in \(\pi\). Detailed balance makes every adjusted marginal exactly \(\pi\). At a mesh size \(\delta\), the expected rejection-weighted displacement under Gibbs is at most \(C\delta^{3/2}\), where \(C\) is finite and independent of sufficiently small \(\delta\). Common-noise proposals are nonexpansive in the small-step regime, so iterating the coupled transition for \(n\) steps gives a joint law of the unadjusted and adjusted endpoints with mean separation at most \(nC\delta^{3/2}\).




Fix a horizon \(T>0\), take \(\delta=T/n\), and realize the unadjusted endpoint using the Brownian mesh increments. The bound becomes


\[\mathbb E\|X_n^{\mathrm{Euler}}-X_n^{\mathrm{adjusted}}\|
\le TC\sqrt\delta\longrightarrow0.\]


The matching-input Euler theorem also gives convergence of the unadjusted endpoint to \(Y_T\) in mean square. For a frequency \(v\), the bounded characteristic test \(x\mapsto e^{i\langle v,x\rangle}\) has Lipschitz constant \(\|v\|\). Its expectations therefore converge both under the vanishing mean-separation coupling and under mean-square Euler convergence. Since the adjusted endpoint has law \(\pi\), the characteristic function of \(Y_T\) equals that of \(\pi\). Uniqueness of characteristic functions proves \(\operatorname{Law}(Y_T)=\pi\); the resulting stationary-law theorem is \cite[Section~3.42.7]{LangevinLean2026ULASC}. At \(T=0\), the conclusion is the initialization assumption.




Hence \(\|g(Y_t)\|_{L^2}\le\sqrt{Ld}\) uniformly in time. The construction supplies both the stationary law and the independent future Brownian increments needed for the one-step comparison.



\section{\texorpdfstring{Integrate the local Euler error with its constants}{Integrate the local Euler error with its constants}}



On an interval of length \(h\), start the stationary process at \(Y_0\). Define the omitted drift integral


\[R_h=\int_0^h\bigl(g(Y_s)-g(Y_0)\bigr)\,ds.\]


The integral equation is exactly \(Y_h=Y_0-hg(Y_0)+\sqrt2B_h-R_h\). Minkowski's inequality and stationarity first give


\[\|Y_s-Y_0\|_{L^2}\le s\sqrt{Ld}+\sqrt{2ds}.\]


Indeed the Gibbs score moment bound (see \cite[Section~3.27.8]{LangevinLean2026ULASC}) and stationarity make the drift integral's norm at most \(\int_0^s\sqrt{Ld}\,du\), while \(\|B_s\|_{L^2}=\sqrt{ds}\). Lipschitzness now controls the integrand defining \(R_h\), so


\[\begin{aligned}
\|R_h\|_{L^2}
&\le L\int_0^h\bigl(s\sqrt{Ld}+\sqrt{2ds}\bigr)\,ds\\
&=L\left(\frac{h^2}{2}\sqrt{Ld}+\frac23h^{3/2}\sqrt{2d}\right).
\end{aligned}\]


If \(Lh\le2\), factor out \(Lh\sqrt{hd}\). The remaining coefficient is at most \(\sqrt2/2+2\sqrt2/3=7\sqrt2/6\le33/20\). Thus


\[\|R_h\|_{L^2}\le\frac{33}{20}Lh\sqrt{hd}.\]


The factor \(h^{3/2}\) is crucial: dividing this local error by a contraction gap of order \(h\) will leave the global \(\sqrt h\) bias.



\section{\texorpdfstring{Compare the discrete and stationary paths}{Compare the discrete and stationary paths}}



Choose any initial coupling of \(X_0\sim\mu_0\) and \(Y_0\sim\pi\) having finite quadratic cost. Extend it with Brownian noise independent of both initial states. Use the increments \((B_{(k+1)h}-B_{kh})/\sqrt h\) as the Gaussian innovations for the numerical chain. Their independence identifies its marginals with \(\mu_k\), regardless of the chosen initial coupling.




Subtracting the integral equation for \(Y\) from the numerical update cancels the common Gaussian increment. With \(T_h(x)=x-hg(x)\),


\[X_{k+1}-Y_{(k+1)h}=T_h(X_k)-T_h(Y_{kh})+R_{k,h}.\]


The sharp deterministic contraction (see \cite[Section~3.10.9]{LangevinLean2026ULASC}) applies pointwise to the first difference, and the local estimate (see \cite[Section~3.11.9]{LangevinLean2026ULASC}) applies to the stationary reference on every interval. Consequently, for \(a_k=\|X_k-Y_{kh}\|_{L^2}\),


\[a_{k+1}\le q a_k+b,\qquad b=\frac{33}{20}Lh\sqrt{hd},\]


where \(q=1-mh\) or \(q=Lh-1\) in the corresponding step regime. Both satisfy \(0\le q<1\). Every \((X_k,Y_{kh})\) is a coupling of \(\mu_k\) with \(\pi\), so its cost bounds the Wasserstein distance. All variables remain square-integrable by the Lipschitz update and the Gaussian second moment.



\section{\texorpdfstring{Sum the recurrence and take the transport infimum}{Sum the recurrence and take the transport infimum}}



Induction on \(k\) gives the exact finite geometric sum


\[a_k\le q^ka_0+b\sum_{j=0}^{k-1}q^j
=q^ka_0+\frac{b(1-q^k)}{1-q}
\le q^ka_0+\frac b{1-q}.\]


The denominator is positive because \(q<1\). This bound holds for every finite-cost initial coupling, and its output laws do not depend on that coupling. Taking the infimum over the initial costs therefore replaces \(a_0\) by \(W_2(\mu_0,\pi)\); an optimal coupling need not be selected.




For the first regime, \(1-q=mh\), hence \(b/(1-q)=(33/20)(L/m)\sqrt{hd}\). For the second regime, \(1-q=2-Lh\), hence \(b/(1-q)=(33/20)Lh\sqrt{hd}/(2-Lh)\). These are exactly the two stated bounds. At \(h=2/(m+L)\), the two contraction factors and the two floors agree.




For fixed positive \(h\), only the initial-error term vanishes as \(k\) grows. To obtain arbitrarily small error one must also reduce \(h\). This explains separately the roles of iteration count and step size.



\begin{thebibliography}{ULA-LSI}
\bibitem[ULA-SC]{LangevinLean2026ULASC}
LangevinLean Project. \emph{Unadjusted Langevin under Strong Convexity}. LL-01, 2026. Draft report, content revision 8d6335dce333. \url{https://odakyildiz.github.io/LangevinLean/docs/pdf/UnadjustedLangevin/FullProof.pdf}.
\end{thebibliography}

\end{document}
